%% ============================================================
%% 片段：s03d-nahm   对应原文 PDF 第 74--82 页
%% 本文件由 gw/convert.py 从 arXiv 源码（hep-th/0612073v2）自动生成：
%%   * 公式/编号/引用/图 已经锁定为原文内容，请勿改动任何数学内容与编号；
%%   * 只需把英文正文译为中文，并在每节开头加“本节重点”、关键处加译注。
%% ============================================================


\subsection{\texorpdfstring{纳姆方程与 $\MH$ 的局部奇点}{纳姆方程与 M_H 的局部奇点}}

%% ==== 原文 §3.8：PDF 第 74 页（印刷第 73 页）====
\label{sec:3.8}

规范场论与纳姆方程（Nahm's equations）可用于在复李群的余伴随轨道（coadjoint orbit）
上得到超 Kähler 度规（hyper-Kähler metric）。原始构造见 \cite{ref34,ref35}，向任意轨道的
推广见 \cite{ref36,ref37}，综述与进一步文献见 \cite{ref63,ref31}。我们关心这些度规，是因为
它们给出了 $\MH(\alpha,\beta,\gamma;p)$ 在局部奇点附近的行为。

构造这些超 Kähler 度规有几条途径。对我们而言，最方便的是考虑穿孔圆盘 $C$ 上的
希钦方程（Hitchin's equations），该圆盘定义为复 $z$ 平面中满足 $|z|\leq 1$ 且 $z\not= 0$
的区域。按惯例记 $z=re^{i\theta}$。在 $C$ 上，我们考虑希钦方程的解，这些解在圆盘的旋转下
不变，并在 $r=0$ 附近具有熟悉的奇异性：
% 式 (3.91)
\setcounter{equation}{90}
\begin{equation}\label{eq:3.91}
  \begin{aligned}A & = \alpha \,d\theta +\dots\\
           \phi & = \beta\,{dr\over r}-\gamma\,d\theta+\dots,\end{aligned}
\end{equation}
设三元组 $(\alpha,\beta,\gamma)$ 是正则（regular）的，即 $G$ 中与该三元组对易的子群恰为
环面 $\TT$。令 $\EUG_{C,p}$ 为旋转不变的规范变换（gauge transformation）$g:C\to G$
所成的群，这些变换在 $|z|=1$ 处等于 1，并在 $z=0$ 处取值于 $\TT$。满足边界条件
\eqref{eq:3.91} 的旋转不变希钦方程解的空间，模去 $\EUG_{C,p}$ 的作用，是一个超 Kähler
流形，我们称它为 ${\mathcal Q}(\alpha,\beta,\gamma)$。

事实上，它可以构造为旋转不变配对 $(A,\phi)$ 构成的空间在 $\EUG_{C,p}$ 作用下的
超 Kähler 商（hyper-Kähler quotient）。这样的配对一般由下式给出：
% 式 (3.92)
\setcounter{equation}{91}
\begin{equation}\label{eq:3.92}
  \begin{aligned}A & = a(r)\,d\theta+h(r){dr\over r}
\\\phi &= b(r){dr\over r}-c(r)\,d\theta\end{aligned}
\end{equation}
其中 $a,b,c$ 与 $h$ 是 $\mathfrak g$-值函数。函数 $a,b,c$ 与通常参数 $\alpha,\beta,\gamma$
以相当显然的方式相联系，但旋转对称性还允许第四个函数 $h$。虽然 $h$ 可以被规范掉，
但目前不这样做更方便。具有这种形式的希钦方程解的空间带有超 Kähler 结构，它就是希钦
方程通常具有的超 Kähler 结构在该情形下的特化。描述它的一种方式是：把函数
$(h,a,b,c)$ 看作从开单位区间出发、取值于四元数（quaternions）$\mathbb{H}\cong\R^4$ 与
$\mathfrak g$ 的张量积之中的映射。超 Kähler 结构来自 $\mathbb{H}$ 上的超 Kähler 结构。
具体而言，在我们将称为 $I$ 的某个复结构（complex structure）中，复变量是 $h-ia$
与 $b+ic$；其余的复结构可由对三元组 $(a,b,c)$ 施加 $SO(3)$ 旋转得到。

若令 $s=-\ln r$ 且 $D/Ds=d/ds+[h,\,\cdot\,]$，则希钦方程变为
% 式 (3.93)
\setcounter{equation}{92}
\begin{equation}\label{eq:3.93}
  \begin{aligned}{Da\over Ds}& = [b,c] \\
                    {Db\over Ds}&=  [c,a] \\
                     {Dc\over Ds}&= [a,b]. \end{aligned}
\end{equation}
若把 $h$ 取为零，它们就成为纳姆方程 \cite{ref64}；这一点局部上可通过规范变换做到。

\begin{yizhu}[易错点]
式~\eqref{eq:3.93}（以及 $h=0$ 时的式~\eqref{eq:3.100}）写的是 $da/ds=[b,c]$ 及其循环
置换。这与文献中常见的写法 $dT_i/ds=\half\epsilon_{ijk}[T_j,T_k]$ 是同一组方程，并不
相差因子：由于 $[\cdot,\cdot]$ 的反对称性，$\half\epsilon_{ijk}[T_j,T_k]$ 中 $(j,k)$ 与
$(k,j)$ 两项相等，求和后恰给出 $[T_j,T_k]$（其中 $(i,j,k)$ 为 $(1,2,3)$ 的循环次序）。
真正需要留意的是生成元的定向与归一化：$t_1,t_2,t_3$ 必须满足 $[t_1,t_2]=t_3$ 及其
循环置换。若改用物理文献中常见的厄米生成元 $T_i=\sigma_i/2$（此时
$[T_1,T_2]=iT_3$），则须取反厄米的 $t_i=-i\sigma_i/2$，否则式~\eqref{eq:3.101} 不满足
式~\eqref{eq:3.100}。（$b+ic=-s^{-1}(t_2+it_3)$ 的幂零性只取决于二次不变量为零，即
$1^2+i^2=0$，与上述定向无关。）此外，本节的 $a,b,c$ 都取值
于 $\mathfrak g$（在四维图像中即 ${\rm ad}(E)$ 的截面），在复结构 $I$ 下把 $b+ic$
解读为希格斯场 $\varphi$ 时尤须记住这一点。
\end{yizhu}

为阐明这些方程的解的模空间 ${\mathcal Q}(\alpha,\beta,\gamma)$ 的性质，首先注意，其中
两个方程的某个线性组合给出
% 式 (3.94)
\setcounter{equation}{93}
\begin{equation}\label{eq:3.94}
  {d\over ds}(b+ic)=-[h-ia,b+ic].
\end{equation}
这表明 $b+ic$ 在 $\mathfrak g_\C$ 中的共轭类（conjugacy class）与 $s$ 无关。我们还有
边界条件 $\lim_{s\to\infty}(b+ic)=\beta+i\gamma$。若 $\sigma=\half(\beta+i\gamma)$ 正则，
则这表明 $\tau=\half(b+ic)$ 处处位于包含 $\sigma$ 的共轭类 $\mathfrak c$ 之中。
特别地，
% 式 (3.95)
\setcounter{equation}{94}
\begin{equation}\label{eq:3.95}
  \tau=\half(b(0)+ic(0))
\end{equation}
就包含在这个共轭类中。$\tau$ 也是规范不变量（因为我们只允许在 $s=0$ 处等于 1 的规范
变换）。于是，把纳姆方程的一个解映到相应的 $\tau$ 值，就得到映射
$\Phi:{\mathcal Q}\to \mathfrak c$，它在复结构 $I$ 下是全纯的。把纳姆方程余下的部分
解释为矩映射（moment map）条件，\cite{ref34} 证明了该映射是同构。

\begin{yizhu}[含糊点拨]
“把纳姆方程余下的部分解释为矩映射条件，从而 $\Phi$ 是同构”一句压缩了本节最要紧的
论证。固定 $s$ 时，旋转不变的希钦方程正是对规范变换群取超 Kähler 商时的矩映射方程；
式~\eqref{eq:3.94} 只给出全纯方向上的演化，而原文所说的“余下部分”即其余两组方程，
它们合起来等价于该商的矩映射条件为零。于是 ${\mathcal Q}$ 作为超 Kähler 流形同构于
$\sigma$ 的伴随轨道，$\Phi$ 就是取解在 $s=0$ 处的值这一映射。同构性以及由此得到的
度规正是轨道上的标准超 Kähler 度规（后者由 $G$ 不变性与轨道辛形式唯一确定），这是
Kronheimer 的定理（文献 [34]，即 P. Kronheimer, \textit{J. London Math. Soc.}
\textbf{42} (1990) 193）的结论，原文未给证明。核对时须注意两点：需要
$(\alpha,\beta,\gamma)$ 的正则性（否则 $\tau$ 只落在从属轨道中，见下一段），以及
$s=0$ 处规范变换取值为 1 的边界条件。可对照本节末尾 $SU(2)$ 的两个极端情形：
$\sigma=\half(\beta+i\gamma)$ 正则时 ${\mathcal Q}\cong G_\C/\TT_\C$（复维数 $2$）；
$\sigma=0$ 时 ${\mathcal Q}\cong T^*(G/\TT)$（$G=SU(2)$ 时即 $T^*S^2$，复维数 $2$）；而在完全非正则的 $\alpha=\beta=\gamma=0$ 处，${\mathcal Q}$ 退化为 $\R^4/\Z_2$ 上的平坦度规（即 $A_1$ 奇点），非零的 $(\alpha,\beta,\gamma)$ 则给出 Eguchi--Hanson 度规，描述 $A_1$ 奇点的形变与消解（详见本节末尾）。
\end{yizhu}

按定义，$\mathfrak c$ 是向量 $\sigma\in\mathfrak g_\C$ 在共轭作用下的轨道（同样也可以
看作对偶空间 $\mathfrak g_\C^\vee$ 中的轨道，因为二次形式 $-\Tr$ 给出 $\mathfrak g_\C$
与 $\mathfrak g_\C^\vee$ 之间 $G_\C$-不变的等同）。它称为伴随（或余伴随）轨道。当
$\sigma$ 正则时，$G_\C$ 中保持 $\sigma$ 不变的子群恰为 $\TT_\C$，即环面 $\TT$ 的复化
（complexification）。因此 $\mathfrak c$ 同构于 $G_\C/\TT_\C$。

若 $\sigma$ 非正则，则 $\tau$ 落在第~\ref{sec:3.3}~节所述的某一从属轨道（affiliated
orbit）$\mathfrak c_\lambda$ 之中。所有可能性都会出现，但对 ${\mathcal Q}$ 中的一般点，
$\tau$ 取值于维数最大的从属轨道，即正则轨道 $\mathfrak c^*$。为了解 $\sigma$ 非正则时
会发生什么，考虑极端情形 $\sigma=0$。此时方程与边界条件允许我们找到 $b$ 与 $c$
恒为零的解。方程退化为 $Da/Ds=0$，于是 $a$ 的共轭类与 $s$ 无关，且等于 $\alpha$ 的
共轭类。这个共轭类必须是正则的，因为我们已经假定 $(\alpha,\beta,\gamma)$ 正则，并且
取了 $\beta=\gamma=0$。现在可以像前面那样推理。$a(0)$ 是规范不变量，并且与 $\alpha$
共轭。此外，模去规范群，$a(0)$ 是 $b=c=0$ 的解的唯一不变量。最后，$a(0)$ 可以是
$\alpha\in\mathfrak g$ 的轨道上的任意元素。对正则的 $\alpha$，该轨道是 $G/\TT$ 的一个
拷贝，这是一个 Kähler 流形，称为旗流形（flag manifold）。为了把 $b$ 与 $c$ 包含进来，
我们注意到，从复结构 $I$ 下的全纯观点看，$b+ic$ 由线性方程 \eqref{eq:3.94} 刻画。
于是 ${\mathcal Q}$ 是 $G/\TT$ 上的全纯向量丛。要使 ${\mathcal Q}$ 成为复辛流形
（并且实际上是超 Kähler 的），该丛必须是余切丛（cotangent bundle）。

于是当 $\beta=\gamma=0$ 而 $\alpha$ 为一般值时，${\mathcal Q}$ 在复结构 $I$ 下就是
$G/T$ 的余切丛：
% 式 (3.96)
\setcounter{equation}{95}
\begin{equation}\label{eq:3.96}
  {\mathcal Q}\cong T^*(G/\TT).
\end{equation}
若形变到 $\beta,\gamma\not=0$，余切丛就形变为 $G/\TT$ 上的仿射丛（affine bundle）。当
$\sigma=\half(\beta+i\gamma)$ 正则时，该仿射丛同构于 $\sigma\in\mathfrak g_\C$ 的
$G_\C$-轨道。

\bigskip\noindent{\it $\mathbb{L}$ 型情形的类比}

将 $\TT$ 换成 $G$ 的任一列维子群（Levi subgroup）$\mathbb{L}$，上述构造都可以重复。

我们要求三元组 $(\alpha,\beta,\gamma)$ 是 $\mathbb{L}$-正则的（$\mathbb{L}$-regular，
即它恰与 $\mathbb{L}$ 对易），并修改 $\EUG_{C,p}$ 的定义，使其由在 $p$ 处取值于
$\mathbb{L}$ 的规范变换组成。记 ${\mathcal Q}_\mathbb{L}(\alpha,\beta,\gamma)$ 为方程 \eqref{eq:3.93} 的解空间，
其边界条件由 $\alpha$、$\beta$、$\gamma$ 设定，模去 $\EUG_{C,p}$ 的作用。

与上面相同的推理表明，若 $\sigma=\half(\beta+i\gamma)$ 是 $\mathbb{L}$-正则的，则
% 式 (3.97)
\setcounter{equation}{96}
\begin{equation}\label{eq:3.97}
  {\mathcal Q}_\mathbb{L}\cong G_\C/\mathbb{L}_\C
\end{equation}
是 $\sigma\in\mathfrak g_\C$ 在 $G_\C$ 的伴随作用下的轨道。在另一极端，若 $\sigma=0$
而 $\alpha$ 为一般值，则
% 式 (3.98)
\setcounter{equation}{97}
\begin{equation}\label{eq:3.98}
  {\mathcal Q}_\mathbb{L}\cong T^*(G/\mathbb{L}).
\end{equation}
一般地，${\mathcal Q}_\mathbb{L}$ 是 $T^*(G/\mathbb{L})$ 的仿射形变。若把 $G$ 换成以
$\mathbb{L}$ 为列维子群的一般半单李群（semi-simple Lie group）$\mathbb{L}'$，则基于
规范群为 $\mathbb{L}'$ 的规范场论的同一构造给出一族超 Kähler 度规，定义在
% 式 (3.99)
\setcounter{equation}{98}
\begin{equation}\label{eq:3.99}
  {\mathcal Q}_{\mathbb{L}',\mathbb{L}}\cong T^*(\mathbb{L}'/\mathbb{L})
\end{equation}
及其仿射形变之上。

\bigskip\noindent{\it $\MH$ 的局部模型}

我们描述过的 $\MH$ 的许多性质，都有涉及超 Kähler 度规 ${\mathcal Q}$ 的局部类比。

例如，$\MH$ 的辛形式（symplectic form）的上同调类（cohomology class）是参数
$\alpha$、$\beta$、$\gamma$ 的线性函数。对 ${\mathcal Q}$ 有类似的线性性，如
\cite{ref34} 的定理 2.6 所述。

再举一例：正如仿射外尔群（affine Weyl group）作用于 $\MH$ 的上同调，普通外尔群
（Weyl group）同样作用于 ${\mathcal Q}$ 的上同调。证明这一点的框架在 \cite{ref31}
第 5 节中描述，其中用到外尔群对参数 $(\alpha,\beta,\gamma)$ 的作用，以及奇点余维为三
（codimension three）这一事实——换言之，正是我们在第~\ref{sec:3.7}~节中用来构造仿射
外尔群在 $\MH$ 的上同调上作用的那些事实。实际上，\cite{ref31} 所考虑的是 \cite{ref35}
中构造的、涉及 Slodowy 切片（Slodowy slice）的一类稍大的超 Kähler 簇。（这些更一般的
簇通过求解与上面相同的方程、但在 $s=0$ 处取不同渐近行为来构造。）对这些簇出现的外尔群
表示称为施普林格表示（Springer representation）。它们可以在一个相对初等的构造
\cite{ref65} 中被几何地理解，而不需要超 Kähler 度规；也许这对 $\MH$ 也有类比。

超 Kähler 流形 ${\mathcal Q}$ 能为 $\MH$ 的许多性质提供局部模型，其根本原因在于：
它们确实描述了带分歧的希格斯丛（ramified Higgs bundle）模空间在局部奇点附近的行为。
例如，如式~\eqref{eq:3.79} 那样，当 $\alpha$ 趋近一个非 $\mathbb{L}$-正则的值
$\bar\alpha$ 时，$\M_{H,\mathbb{L}}(\alpha,0,0;p)$ 的局部行为由
$T^*(\mathbb{L}'/\mathbb{L})$ 建模，其中 $\mathbb{L}'$ 是某个真正包含 $\mathbb{L}$ 的
列维群。为了得到该情形的局部模型，我们需要 $T^*(\mathbb{L}'/\mathbb{L})$ 上一族合适的
超 Kähler 度规（以及 $\beta,\gamma\not=0$ 时的仿射形变）。由上述构造正好可以得到它，
如式~\eqref{eq:3.99} 所指出的，只要取规范群为 $\mathbb{L}'$。用 $\mathbb{L}'$ 规范
理论而不是 $G$ 规范场论来构造局部模型有一个简单的道理：奇点处的边界条件（在此情形由
三元组 $(\bar\alpha,0,0)$ 定义）只在 $\mathbb{L}'$ 下不变，而不在 $G$ 下不变，因此以
$\mathbb{L}'$ 为规范对称性的理论就足以描述该奇点。

\bigskip\noindent{\it $SU(2)$ 的详细分析}

为了更充分地理解为什么簇 ${\mathcal Q}$ 能很好地描述 $\MH$ 的局部奇点，我们将更仔细地
考察 ${\mathcal Q}$ 在非正则点附近的行为。为简单起见，我们以 $G=SU(2)$ 为例来说明。

对 $G=SU(2)$，三元组 $(\alpha,\beta,\gamma)$ 的唯一非正则值是
$\alpha=\beta=\gamma=0$。在 $\alpha=\beta=\gamma=0$ 处，${\mathcal Q}$ 描述圆盘
$|z|\leq 1$ 上旋转不变、且在 $r=0$ 处比 $1/r$ 更不奇异的希钦方程解。

等价地，取规范 $h=0$，我们需要纳姆方程
% 式 (3.100)
\setcounter{equation}{99}
\begin{equation}\label{eq:3.100}
  \begin{aligned}{da\over ds}&=[b,c]\\
                      {db\over ds}&=[c,a]\\
                      {dc\over ds}&=[a,b],\end{aligned}
\end{equation}
在半直线 $[0,\infty)$ 上的解，且当 $s\to \infty$ 时 $a,b,c\to 0$。

显然，一个这样的解是 $a=b=c=0$。另一个简单解（它是 \cite{ref35} 的出发点）是
% 式 (3.101)
\setcounter{equation}{100}
\begin{equation}\label{eq:3.101}
  \begin{aligned}a&=-{1\over s}t_1\\
                    b&=-{1\over s}t_2\\
                    c&=-{1\over s}t_3,\end{aligned}
\end{equation}
其中 $t_1,t_2,t_3$ 是 ${\mathfrak sl}(2,\C)$ 中满足 $[t_1,t_2]=t_3$ 及其循环置换的
固定元素。

在复结构 $I$ 下，这个解描述一个希格斯丛 $(E,\varphi)$，其中 $\varphi$ 具有留数为幂零
（nilpotent）的极点。（注意 $b+ic=-s^{-1}(t_2+it_3)$ 对所有 $s$ 确实都是幂零的。）在
复结构 $J$ 下，它描述一个平坦丛（flat bundle），其绕 $r=0$ 处的单值性（monodromy）是
幂幺（unipotent）的但不等于 1。这些结果正是从第~\ref{sec:3.3}~节可以预期的。

一个小的推广是引入正常数 $f$ 并取
% 式 (3.102)
\setcounter{equation}{101}
\begin{equation}\label{eq:3.102}
  \begin{aligned}a&= -{1\over s+f^{-1}}t_1\\
                    b&=-{1\over s+f^{-1}}t_2\\
                    c&=-{1\over s+f^{-1}}t_3,\end{aligned}
\end{equation}
我们之所以用这种特定方式参数化解，是因为 \eqref{eq:3.102} 在 $f\to 0$ 时确实有极限，
即平凡解 $a=b=c=0$。因此 $f$ 取值于 $\R_{\geq 0}=[0,\infty)$。在我们的应用中，我们
关心的是小 $f$。我们还可以选取元素 $R\in SO(3)$，把 \eqref{eq:3.102} 推广为
% 式 (3.103)
\setcounter{equation}{102}
\begin{equation}\label{eq:3.103}
  \begin{aligned}a&= -{1\over s+f^{-1}}Rt_1R^{-1}\\
                    b&=-{1\over s+f^{-1}}Rt_2R^{-1}\\
                    c&=-{1\over s+f^{-1}}Rt_3R^{-1}.\end{aligned}
\end{equation}
于是这一族的参数空间是 $\R_{\geq 0}\times SO(3)=\R^4/\Z_2$，其中 $\Z_2$ 通过作用于
$\R^4$ 全部四个坐标上的反射来作用。这就是 $A_1$ 奇点。特别地，在这个由纳姆
方程或希钦方程的解构成的族 ${\mathcal Q}$ 上，通过在圆盘 $r\leq 1$ 上积分场 $(A,\phi)$
变分的 ${\mathfrak{L}}^2$ 范数得到的自然度规，是 $\R^4/\Z_2$ 上的平坦度规；所有模式
都是平方可积的，这要归功于因子 $1/s=1/(-\ln r)$，而 $f=0$ 处于有限距离处。对非零的
$\alpha,\beta,\gamma$，我们得到的则是 Eguchi--Hansen 度规，它描述 $A_1$ 奇点的形变与
消解（resolution）。

\begin{yizhu}[原文疑误]
原文此处写作“Eguchi-Hansen 度规”，通行拼写为 Eguchi--Hanson 度规（T. Eguchi 与
A. J. Hanson, 1978）；作者在原文其他处（如第 3.6 节讨论 Calabi 度规的段落）也写作
“Eguchi and Hansen”，可见并非源码转换之误，而系原稿用字。译文一律保留原拼写，以免与
原文其他处不一致；出版整理时建议统一改为 Eguchi--Hanson。
\end{yizhu}

现在可以解释为什么纳姆方程与 Kronheimer 构造能为 $\MH$ 的局部奇点提供良好的模型。
${\mathcal Q}$ 的奇点出现在 $f=0$ 处，对应于平凡解 $a=b=c=0$。${\mathcal Q}$ 中奇点
附近的点对应于 $f$ 非常小的解。当纳姆方程嵌入希钦方程时，多出一个 $1/r$ 因子，这在
式~\eqref{eq:3.92} 中已可见到。这意味着对小的 $f$，除 $r\lesssim f$ 之外所有场都非常
小。若把 $\alpha,\beta,\gamma$ 扰动为非零但为 $f$ 阶的量，情况仍然如此。因此，我们可以
从任意无分歧的希格斯丛 $(E,\varphi)$ 出发，把上面的解族“嫁接”进去，从而得到一族带
分歧的希格斯丛。对 $r>>f$ 我们用原来的解，而对 $r\lesssim f$ 用精确解 \eqref{eq:3.103}。
这样得到的解族在令 $f=0$ 时获得一个 $A_1$ 奇点，此时“新”解退化为旧解。

当然，嫁接操作给出的并不是精确解。由嫁接得到的场集合必须经过修正才能成为精确解，但当
$f$ 小时所需的修正也小。这有点类似于 \cite{ref66} 中在四维流形上构造杨–米尔斯
（Yang-Mills）瞬子（instanton）解的做法：把一个主要支撑在 $\R^4$ 中非常小区域上的精确
解嫁接进去。然而，与瞬子模空间不同，希格斯丛的形变理论是无障碍的（unobstructed）（只要
远离可约希格斯丛），因此不存在 \cite{ref66} 中所述的那种可能阻碍把嫁接场形变为精确解的
拓扑条件。

\begin{yizhu}[物理图像补充]
这里的“嫁接”图像值得记住：$r\lesssim f$ 的区域由纳姆方程的解~\eqref{eq:3.103} 主导，
$r\gg f$ 的区域由原有的无分歧希格斯丛 $(E,\varphi)$ 主导，两者在 $r\sim f$ 处衔接。
由于希格斯丛的形变理论在远离可约丛时无障碍，衔接后的场可以逐次微扰为精确解；这正是
$A_1$ 奇点（及其在 $f\neq 0$ 时的消解）能作为 $\MH$ 局部模型的原因。与 Taubes 在四维
流形上拼接瞬子的构造相比，这里无需满足任何可能阻碍形变的拓扑条件。$f$ 控制奇点邻域的
尺度：$f\to 0$ 时邻域收缩、解退化为无分歧解；把 $\alpha,\beta,\gamma$ 扰动到 $f$ 阶时，
该邻域的度规变为 Eguchi--Hansen 型，奇点被形变并消解。另外，“奇异轨迹恰为 $\MH$”一语
应理解为：带分歧模空间在整个无分歧模空间 $\MH$（而不只是某一点）上方都出现 $A_1$ 型
奇点。
\end{yizhu}

这样，我们就得到了一个更精确的方式来理解第~\ref{sec:3.6}~节所论证的内容：当
$\alpha,\beta,\gamma\to 0$ 时，$\MH(\alpha,\beta,\gamma;p)$ 出现 $A_1$ 奇点，其奇异
轨迹恰为 $\MH$，即无分歧希格斯丛的模空间。这一具体结果当然是 $G=SU(2)$ 特有的。对更高
秩的 $G$，正如我们所论证的，类似的构造一般会导致更复杂的奇点，其消解为对各种
$\mathbb{L}'$、$\mathbb{L}$ 的 $T^*(\mathbb{L}'/\mathbb{L})$。

\subsection{希钦纤维化}
%% ==== 原文 §3.9：PDF 第 80 页（印刷第 79 页）====
\label{sec:3.9}

下面我们在带分歧的希格斯丛的背景下描述希钦纤维化（Hitchin fibration）与完全可积性
（complete integrability）。所有这些内容在文献中都有更详细的讨论；例如见
\cite{ref67,ref68}。

我们从 $SL(2,\C)$ 的例子开始，见第~\ref{sec:3.3}~节。在共轭意义下，$\varphi$ 在
$\sigma\not=0$ 的抛物点（parabolic point）附近的局部行为为
% 式 (3.104)
\setcounter{equation}{103}
\begin{equation}\label{eq:3.104}
  \varphi={dz\over z}\sigma(1+{\mathcal O}(z)).
\end{equation}
由此可知
% 式 (3.105)
\setcounter{equation}{104}
\begin{equation}\label{eq:3.105}
  \Tr\,\varphi^2=\Tr\,\sigma^2 \left({dz\over z}\right)^2(1+\CO(z)).
\end{equation}
这个论断实际上对所有 $\sigma$（无论其是否为零）都一致成立。确实，当 $\sigma=0$ 时，
$\varphi$ 的极点部分（polar part）是幂零的，但在一般解中 $\varphi$ 还有正则项。把
$z=0$ 附近的全纯结构平凡化之后，在共轭意义下有
% 式 (3.106)
\setcounter{equation}{105}
\begin{equation}\label{eq:3.106}
  \varphi={dz\over z}\left(\left(\begin{matrix}0&1\\ 0&0\end{matrix}\right)+az+bz^2+\dots\right),
\end{equation}
其中 $a,b\in {\mathfrak sl}(2,\C)$。因此二次微分（quadratic differential）
$\Tr\,\varphi^2$ 在 $z=0$ 处可以有单极点，但没有二阶极点，这表明 \eqref{eq:3.105}
对 $\sigma=0$ 也成立。

这一点之所以重要，是因为二次微分 $\Tr\,\varphi^2$ 是 $\MH$ 作为复结构 $I$ 下的复流形
完全可积的关键。在没有分歧的情况下，对 $G_\C=SL(2,\C)$，$\MH$ 的（复）维数等于
$6g_C-6$，其中 $g_C$ 是 $C$ 的亏格（genus）。为建立完全可积性，需要 $3g_C-3$ 个对易的
哈密顿量（commuting Hamiltonian）。它们正是 $\Tr\,\varphi^2$ 的各分量，而
$\Tr\,\varphi^2$ 取值于 $C$ 上二次微分构成的空间。该空间的维数是 $3g_C-3$。（关于
希钦系统完全可积性的极简说明，见 \cite{ref8} 第 4.3 节。）

下面在带分歧的情形做类似的计算。为叙述简单，设只有一个分歧点 $p$。希格斯丛是一个
配对 $(E,\varphi)$，其中 $E$ 是 $C$ 上的 $SL(2,\C)$-丛，而
$\varphi\in H^0(C,K_C\otimes {\rm ad}(E)\otimes \CO(p))$。这里我们纳入因子 $\CO(p)$
（其截面是在 $p$ 处允许有单极点的函数构成的丛），因为 $\varphi$ 允许在 $p$ 处有极点。
确定丛 $E$ 所需的参数个数是 $3g_C-3$。由黎曼–罗赫定理（Riemann-Roch），
$H^0(C,K_C\otimes {\rm ad}(E)\otimes \CO(p))$ 的维数是 $3g_C$。（确实，
$H^0(C,K_C\otimes{\rm ad}(E))$ 是稳定丛模空间的余切空间，维数为 $3g_C-3$；与 $\CO(p)$
的张量积使维数增加 3，因为 $K_C\otimes {\rm ad}(E)$ 的秩为 3。）然而 $\varphi$ 服从一个
约束 \eqref{eq:3.106}。因此 $\varphi$ 的选择依赖于 $3g_C-1$ 个参数，而
$\MH(\alpha,\beta,\gamma;p)$ 的维数是 $(3g_C-3)+(3g_C-1)=6g_C-4$。于是为建立完全可积性
需要 $3g_C-2$ 个对易哈密顿量。

这些正是 $\Tr\,\varphi^2$ 的各分量。一般地，$\Tr\,\varphi^2$ 是一个带双极点的二次
微分。这样的二次微分构成的空间维数为 $3g_C-1$，但其中一个参数由 \eqref{eq:3.105}
确定，因此 $\Tr\,\varphi^2$ 取值于一个 $3g_C-2$ 维的空间。这就给出了
$\MH(\alpha,\beta,\gamma;p)$ 完全可积所需的 $3g_C-2$ 个参数。

关于 $\MH$ 的希钦纤维化、完全可积性等等的所有一般性论述，对带分歧的希格斯丛都有自然
的类比。我们记 $\EUBB$ 为 $p$ 处有双极点、且满足 \eqref{eq:3.105} 的二次微分构成的
空间。希钦纤维化是映射 $\pi:\MH(\alpha,\beta,\gamma;p)\to\EUBB$，它在复结构 $I$ 下
全纯。它由把希格斯丛 $(E,\varphi)$ 映到 $\EUBB$ 中由 $\Tr\,\varphi^2$ 定义的点而
给出。$\EUBB$ 上的函数即对易哈密顿量，而纤维是复拉格朗日子流形（complex Lagrangian
submanifold）（也就是说，从全纯二形式 $\Omega_I$ 的角度看它们是拉格朗日的）。一般
纤维是阿贝尔簇（abelian variety）。此外，若在多个点带分歧（而不只一个点），这一切
仍然成立。

对任意 $G$，其类比如下。设 $r$ 是 $G$ 的秩。李代数 $\mathfrak g$ 上的不变多项式
（invariant polynomial）环由 $r$ 个多项式 ${\mathcal P}_i$ 自由生成，它们对某些整数
$d_i$ 是 $d_i$ 次齐次的。这些整数满足
% 式 (3.107)
\setcounter{equation}{106}
\begin{equation}\label{eq:3.107}
  \sum_{i=1}^r(2d_i-1)={\rm dim}(G).
\end{equation}
代替二次微分 $\Tr\,\varphi^2$，我们考虑其类比 ${\mathcal P}_i(\varphi)$，它们是
$K_C^{d_i}$ 的全纯截面。它们满足
% 式 (3.108)
\setcounter{equation}{107}
\begin{equation}\label{eq:3.108}
  {\mathcal P}_i(\varphi)=
{\mathcal P}_i(\sigma)\left({dz\over z}\right)^{d_i}(1+{\mathcal O}(z)).
\end{equation}
这些条件对所有 $\sigma$ 成立，并且，对 $\mathbb{L}=\TT$，微分 ${\mathcal P}_i(\varphi)$
除此之外在 $z=0$ 处可以有 $d_i$ 阶的任意极点。对其他 $\mathbb{L}$，极点还须满足额外
条件。这里我们将集中讨论 $\mathbb{L}=\TT$ 的情形。$\EUBB$ 的定义是：$\EUBB$ 中的点
是一族 $d_i$-微分，$i=1,\dots, r$，它们在分歧点之外全纯，并在这种点附近的行为如
\eqref{eq:3.108} 所示。$\EUBB$ 的维数（对只有一个分歧点的情形）是
$(g_C-1){\rm dim}(G)+\sum_{i=1}^r(d_i-1)$，这可借助 \eqref{eq:3.107} 证明。

$\MH$ 的维数可用我们处理 $SU(2)$ 时的办法计算。$G$-丛 $E$ 的选择依赖于
$(g_C-1){\rm dim}(G)$ 个参数。$H^0(C,K_C\otimes {\rm ad}(E)\otimes\CO(p))$ 的维数是
$g_C{\rm dim}(G)$，但 $\varphi$ 受 $r$ 个约束 \eqref{eq:3.108} 的限制，因此它依赖于
$g_C{\rm dim}(G)-r$ 个参数。于是 $\MH(\alpha,\beta,\gamma;p)$ 的维数为
% 式 (3.109)
\setcounter{equation}{108}
\begin{equation}\label{eq:3.109}
  \begin{aligned}{\rm dim}(\M_H(\alpha,\beta,\gamma;p))
=&\,(2g_C-1)\,{\rm dim}(G)-r =(2g_C-2)\,{\rm dim}(G)+\sum_{i=1}^r(2d_i-2)\\=&
\,{\rm dim}(\MH)+{\rm dim}(G)-r.\end{aligned}
\end{equation}
为建立完全可积性所需的对易哈密顿量的数目因此是
$(g_C-1){\rm dim}(G)+\sum_{i=1}^r(d_i-1)$。这恰好是 $\EUBB$ 的维数。希钦纤维化
$\pi:\MH(\alpha,\beta,\gamma;p)\to \EUBB$ 由把配对 $(E,\varphi)$ 映到 $\EUBB$ 中由
$({\mathcal P}_1(\varphi),\dots,{\mathcal P}_r(\varphi))$ 定义的点而给出，它具有我们为
$G=SL(2,\C)$ 描述过的全部关键性质。

\begin{yizhu}[易错点]
（1）引用提示：上文“$\varphi$ 服从一个约束 \eqref{eq:3.106}”所引的 (3.106) 是 $\sigma=0$
情形的局部形式——被固定的恰是它的领头项（这正是那个“约束”），而 $a,b,\dots$ 等高阶系数
仍然自由；本段末尾与下一段两次引用的是与之等价的 (3.105)（即 $\Tr\,\varphi^2$ 在 $p$ 处的
双极点系数固定），对一般的 $\sigma$ 更准确的一般形式见 (3.104)。两处指同一约束，故照原文
保留 (3.106)，建议出版时核对 PDF 第 81 页。
（2）基空间维数易错：$p$ 处有双极点的二次微分空间维数为 $3g_C-1$，但式~\eqref{eq:3.105}
固定了其中一个参数，故 $\EUBB$ 的维数是 $3g_C-2$，不可与 $3g_C-1$ 混淆；对一般的 $G$，
基空间维数为 $(g_C-1)\dim(G)+\sum_{i=1}^r(d_i-1)$（用式~\eqref{eq:3.107} 化简即得），
它恰等于对易哈密顿量的个数，也是完全可积性的标志。
（3）判别式与奇异纤维：$\EUBB$ 中使谱曲线（spectral curve，即特征多项式）出现重根的
点构成判别式（discriminant）轨迹，其上方的希钦纤维是退化的奇异纤维；一般纤维则是
阿贝尔簇。第 3.6--3.8 节讨论的、当参数趋于非正则值时出现的 $\MH$ 局部奇点，与判别式
轨迹附近的几何密切相关：判别式上方纤维退化，其局部模型正是
$T^*(\mathbb{L}'/\mathbb{L})$；对 $G=SL(2,\C)$，判别式就是 $\Tr\,\varphi^2$ 有重零点的
条件。
\end{yizhu}

式~\eqref{eq:3.109} 给出的 $\MH(\alpha,\beta,\gamma;p)$ 的（复）维数公式，当然与在
适当极限下存在一般纤维化 \eqref{eq:3.79} 相容。对更一般的列维子群 $\mathbb{L}$，对此
有自然的修正，即对约束 \eqref{eq:3.108} 作适当的扩充。
